% swim-times-data-gap.tex
% Technical report: reachability of swim date, meet, and time standard
% Compiled with: pdflatex swim-times-data-gap.tex

\documentclass[12pt]{article}

\usepackage{geometry}
\geometry{letterpaper, margin=1in}

\usepackage{hyperref}
\usepackage{booktabs}
\usepackage{array}
\usepackage{longtable}
\usepackage{enumitem}
\usepackage{titlesec}
\usepackage{fancyhdr}
\usepackage{xcolor}
\usepackage{listings}

\definecolor{castblue}{RGB}{0,48,135}
\definecolor{lightgray}{gray}{0.95}
\definecolor{okgreen}{RGB}{0,120,0}
\definecolor{badred}{RGB}{170,0,0}

\lstset{
  basicstyle=\ttfamily\small,
  backgroundcolor=\color{lightgray},
  frame=single,
  framesep=4pt,
  rulecolor=\color{lightgray},
  breaklines=true,
  showstringspaces=false,
  language=,
}

\pagestyle{fancy}
\fancyhf{}
\renewcommand{\headrulewidth}{0.4pt}
\fancyhead[L]{\small College Area Swim Team}
\fancyhead[R]{\small Data-Gap Report: swim-times}
\fancyfoot[C]{\thepage}

\titleformat{\section}{\large\bfseries\color{castblue}}{\thesection}{0.5em}{}[\titlerule]
\titleformat{\subsection}{\normalsize\bfseries}{\thesubsection}{0.5em}{}

\newcommand{\code}[1]{\texttt{#1}}
\newcommand{\yes}{\textcolor{okgreen}{\textbf{200 OK}}}
\newcommand{\no}{\textcolor{badred}{\textbf{403}}}

\begin{document}

% -----------------------------------------------------------------
\begin{titlepage}
  \centering
  \vspace*{2cm}
  {\Huge\bfseries\color{castblue} Technical Report\\[0.5em]}
  {\Large Retrieving Swim Date, Meet, and Time Standard\\[0.3em]
   from the USA~Swimming Data Hub}\\[2em]
  {\large College Area Swim Team (CAST)}\\[0.5em]
  {\normalsize Prepared for: J.~Evans (project owner)}\\[0.5em]
  {\normalsize Prepared by: CAST Technology Working Group}\\[2em]
  \begin{tabular}{ll}
    Document identifier: & CAST-TR-002                                 \\
    Revision:            & 1.1                                         \\
    Date:                & July 13, 2026                               \\
    Status:              & Released                                    \\
    Subject program:     & \code{swim2/swim-times.w}                    \\
  \end{tabular}
  \vfill
  {\small\textit{All endpoint behaviour reported here was probed live against
  the production API on 13~July~2026 using the program's own anonymous
  authentication headers.}}
\end{titlepage}

\tableofcontents
\clearpage

% =================================================================
\section{Executive Summary}
% =================================================================

The \code{swim-times} program fetches each swimmer's best time per event
from the USA~Swimming data hub, but leaves three columns empty: the
\textbf{swim date}, the \textbf{meet name}, and the \textbf{motivational /
age-group time standard} (B, BB, A, \dots). This report determines whether
those fields can be obtained and, where they cannot, explains why and lays
out alternatives.

\medskip
\noindent\textbf{Verdict.}
\begin{itemize}[leftmargin=1.4em]
  \item The three fields \emph{do} exist in the API and are even returned
        \emph{anonymously} by one endpoint
        (\code{AllTimes/SwimTime/\{swimTimeId\}}), which yields
        \code{swimDate}, \code{meetName}, and \code{timeStandard} together.
  \item \textbf{They cannot be reached anonymously for a given swimmer's
        best times}, because the only anonymous per-swimmer endpoint returns
        a \code{swimTimeRecognitionId} that lives in a \emph{different id
        space} than the \code{swimTimeId} the detail endpoint requires, and
        every endpoint that would enumerate a member's \code{swimTimeId}s
        (full-times, meets, per-member standards) rejects the anonymous
        caller with HTTP~403.
  \item \textbf{The time-standard column is fully solvable offline} with data
        already in the repository, independent of the API.
  \item \textbf{Date and meet are obtainable only with an authenticated
        USA~Swimming session}, which is the sole complete fix.
\end{itemize}

% =================================================================
\section{Background: What Is Missing and Why}
% =================================================================

The anonymous best-times feed the program relies on
(\code{GetBestTimesForMember}) returns exactly six fields per event:

\begin{lstlisting}
{ "swimTimeRecognitionId": 170567980, "strokeName": "Freestyle",
  "strokeAbbreviation": "FR", "distance": 50,
  "courseCode": "SCY", "swimTime": "22.64r" }
\end{lstlisting}

There is no date, no meet, and no standard in this record. The program
therefore prints those columns empty, and the automated test suite records
the corresponding failures (\code{REQ-F-42}, \code{REQ-F-14/15/16} for the
Ledecky age-window filters, and \code{REQ-F-73} for de-duplication, whose
\code{UNIQUE} key includes the absent \code{swim\_date}).

% =================================================================
\section{Investigation Method}
% =================================================================

The production times API at \code{times-api.usaswimming.org} authenticates
anonymous callers with three headers --- \code{Usas-Sub-Id: Anonymous},
\code{AppName: DataHub}, and a format-validated \code{Device-Id} --- exactly
as the program does. Using those headers, every endpoint in the published
OpenAPI description (\code{/swagger/v1/swagger.json}, 122 paths) that could
plausibly carry date, meet, or standard was exercised with representative
parameters. Katie Ledecky (member \code{6CD35348E5824C}) was used as a
stable, public test subject.

% =================================================================
\section{Findings}
% =================================================================

\subsection{Endpoint reachability (anonymous)}

\begin{center}
\begin{longtable}{@{}p{6.9cm} l p{4.6cm}@{}}
  \toprule
  \textbf{Endpoint} & \textbf{Status} & \textbf{Carries date/meet/std?} \\
  \midrule\endhead
  \code{GetMembersForFilters} (search)      & \yes & no (identity only) \\
  \code{GetMember/\{memberId\}}             & \yes & no \\
  \code{GetBestTimesForMember/\{memberId\}} & \yes & \textbf{no} --- time only \\
  \code{AllTimes/SwimTime/\{swimTimeId\}}   & \yes & \textbf{YES --- all three} \\
  \code{GetAllEvents}, \code{GetCompetitionGenders} & \yes & reference data \\
  \code{GetAllTimeStandards}, \code{GetSeasons}     & \yes & reference data \\
  \midrule
  \code{GetAllTimesForFilters} (full history) & \no & would, but blocked \\
  \code{Time/search}, \code{Time/TopTimes/search} & \no & would, but blocked \\
  \code{GetSwimmerMeets/\{memberId\}}       & \no & would, but blocked \\
  \code{GetSwimmerCourses/\{memberId\}}     & \no & blocked \\
  \code{GetSwimmerEvents/\{memberId\}}      & \no & blocked \\
  \code{GetSwimmerTimeStandards/\{memberId\}} & \no & blocked \\
  \code{Time/Recognition/\{id\}}, \code{Time/\{id\}} & \no & blocked \\
  \code{BiqQuery/GetBestTimesForMember}     & \no & blocked \\
  \bottomrule
\end{longtable}
\end{center}

\subsection{The complete-detail endpoint works anonymously}

\code{AllTimes/SwimTime/\{swimTimeId\}} returns everything the program is
missing, without any credential:

\begin{lstlisting}
GET /swims/TimesSearch/AllTimes/SwimTime/170567980   ->  200 OK
{ "meetName": "2019 OZ CFAC MLK Invitational",
  "eventCode": "100 BK SCY", "swimDate": "Jan 19, 2019",
  "fullName": "Brendan Worster", "swimTime": "1:01.61",
  "timeStandard": "A", "timeDrop": -4.02,
  "finishPosition": 2, "powerPoints": 591 }
\end{lstlisting}

\subsection{The core obstacle: two different id spaces}

The trap is subtle. Ledecky's best 50~FR~SCY record above carries
\code{swimTimeRecognitionId = 170567980}. Feeding that number to the detail
endpoint returns \emph{a different swimmer entirely} (Brendan Worster's
100~BK~SCY, shown above). The \code{swimTimeRecognitionId} from the
best-times feed and the \code{swimTimeId} the detail endpoint consumes are
\textbf{distinct identifier spaces that merely overlap numerically.} There
is no anonymous endpoint that converts one to the other, and no anonymous
endpoint that lists a member's \code{swimTimeId}s. The endpoints that would
(\code{GetAllTimesForFilters}, \code{Time/search}, \code{GetSwimmerMeets})
all answer \no\ regardless of the request body --- the rejection is on the
\code{Anonymous} subject, not on parameters.

\medskip
\noindent\textbf{Conclusion:} the detail endpoint is a locked room with the
data inside it; the anonymous tier hands us the wrong key and withholds the
directory that would tell us which key to cut.

% =================================================================
\section{Options and Alternatives}
% =================================================================

\subsection{Option A --- Authenticated access (the only complete fix)}

Log in with a USA~Swimming account through \code{security-api} and use the
real subject id instead of \code{Anonymous}. With a valid sub,
\code{GetAllTimesForFilters} returns the full swim history --- each row
carrying \code{swimTimeId}, \code{swimDate}, \code{meetName}, and
\code{timeStandard} --- which closes all three gaps at once and also revives
the Ledecky age-window filters (\code{REQ-F-14/15/16}) and de-duplication
(\code{REQ-F-73}).

\begin{itemize}[leftmargin=1.4em]
  \item \textbf{Effort:} medium. Add a login step (\code{Account/Login} ->
        token), send the token / real \code{Usas-Sub-Id}, and switch the
        fetch path from best-times to \code{GetAllTimesForFilters}.
  \item \textbf{Cost to you:} your USA~Swimming credentials, stored/entered
        securely (e.g.\ an environment variable, never in the source).
  \item \textbf{Caveat:} confirm that automated retrieval from
        logged-in endpoints is within USA~Swimming's acceptable-use terms
        before deploying it.
\end{itemize}

\subsection{Option B --- Compute the standard locally (do this regardless)}

The \emph{time standard} column needs no API at all. The repository already
contains \code{motivational-time-standards.pdf} (B/BB/A/\dots) and
\code{sdi-agc-time-standards.pdf}. Given a best time plus the swimmer's age
and gender and the event, the attained standard is a deterministic table
lookup. This fully closes the standard column offline and would make
\code{REQ-F-32}'s standard field meaningful even in anonymous mode.

\begin{itemize}[leftmargin=1.4em]
  \item \textbf{Effort:} low--medium. Encode the two standards tables (they
        already exist as PDFs) and compare.
  \item \textbf{Limitation:} solves standard only; date and meet still
        require Option~A.
\end{itemize}

\subsection{Option C --- Enrich by swimTimeId when one is available}

Because \code{AllTimes/SwimTime/\{swimTimeId\}} is anonymous and complete,
\emph{any} route that yields a real \code{swimTimeId} makes date, meet, and
standard a single extra GET away. The practical routes to a
\code{swimTimeId} are Option~A (login) or scraping (Option~E) --- so this is
an enrichment step layered on another option, not a standalone solution.

\subsection{Option D --- Public leaderboard (partial, ranked swims only)}

\code{GetTopTimesLeaderBoard} is public in principle and its rows include a
\code{swimTimeId}. In testing it required an exact combination of season,
LSC, and zone parameters and otherwise returned ``no records,'' and it only
covers \emph{top-ranked} swims per event. It is therefore unusable for
arbitrary age-group swimmers and is not recommended as the primary path.

\subsection{Option E --- Read the website / manual lookup}

The data hub website renders date, meet, and standard for anonymous
visitors. A person can read them directly (the manual process the program
was built to replace), or a browser-automation scraper could drive the
JavaScript single-page app. Scraping was rejected earlier as brittle
(JS-rendered content, no stable contract) and is a last resort.

% =================================================================
\section{Recommendation}
% =================================================================

\textbf{Implementation status (Rev.~1.1).} Option~B has since been
implemented for the motivational standard: the official standards are loaded
into a \code{motivational\_standard} Postgres table (seeded from the PDF by
\code{tools/parse-motivational.py}), and the program computes each best
time's standard with one \code{libpq} query, keyed on the swimmer's gender
and the age group derived from the live member search. The standard column
is now populated; date and meet remain empty pending the Option~A decision.

\begin{enumerate}[leftmargin=1.6em]
  \item \textbf{Option~B --- done.} The motivational standard is computed
        locally from a Postgres table seeded from the official standards,
        closing one of the three columns with no credentials or network risk.
  \item \textbf{Decide on Option~A for date and meet.} If those two columns
        matter, authenticated access is the only complete route; it also
        restores the age-window and de-duplication requirements. If they do
        not, leave them empty and mark \code{REQ-F-42} and \code{REQ-F-14/15/16}
        as \emph{out of scope for the anonymous build}, which the SRS and SDD
        already document as a known limitation.
  \item \textbf{Do not pursue Options~D or~E} except as a fallback; neither
        is a reliable general solution.
\end{enumerate}

\medskip
\noindent In short: the standard is a quick offline win; the date and meet
are gated behind a USA~Swimming login and require a product decision about
whether to authenticate.

\end{document}
